\documentclass[10pt,a4paper]{amsart}
\usepackage[margin=19mm]{geometry}
\usepackage{amsmath,amssymb,mathtools}
\usepackage[scr=rsfs,cal=euler]{mathalfa}
\usepackage{xfrac}
\usepackage[unicode,hidelinks,bookmarks=false]{hyperref}
\newcommand{\ie}{i.e.\ }
\newcommand{\cf}{cf.\ }
\newcommand{\resp}{respectively\ }
\newcommand{\Subsection}{\subsection}
\providecommand{\nobreakdash}{\nobreak\hbox{-}}
\setlength{\emergencystretch}{2em}
\multlinegap0pt
\usepackage{bm}
\usepackage{tikz}
\usepackage{xcolor}
\usepackage{float}
\usepackage{amssymb}
\usepackage{mathtools}
\usepackage[shortlabels]{enumitem}
\usepackage{tcolorbox}
\newtheorem{lemma}{Lemma}
\newcommand{\JSCheck}[1]{#1}
\title{Asymmetric Scaling Laws from Sparse Features}
\author{John Sous, Michael Winer}
\date{Source: arXiv:2605.23591v1 (CC BY 4.0)}
\begin{document}
\maketitle

\noindent\emph{Source and licence.} Asymmetric Scaling Laws from Sparse Features. Version arXiv:2605.23591v1 (CC BY 4.0), distributed by arXiv under the Creative Commons Attribution 4.0 International licence, http://creativecommons.org/licenses/by/4.0/, as recorded in the arXiv licence metadata for this exact version and shown on the abstract page https://arxiv.org/abs/2605.23591v1. Full licence notice: \texttt{licenses/arxiv-2605.23591-CC-BY-4.0.txt}.\par\smallskip

\noindent\textit{Editorial scope.} Counted unit: the article's lemma lem:learnable-coordinates (source0.tex lines 357--370). The article's complete original proof of it is delivered with it (source0.tex lines 976--1007, one \texttt{proof} environment of 32 lines, which the article places away from the statement). No mathematics is written, repaired or re-derived: the statement and the proof are the article's own lines, in the article's own notation, and no retained line is substituted or re-flowed.  Retained uncounted as the source-local context the proof needs: lines 293--296; lines 378--386.  Nothing was omitted from any retained span, so the package carries no omission record.  No retained line contains a citation, so the wrapper carries no bibliography and no bibliography entry.  No retained line contains a figure environment, an image-inclusion command or drawing code, so no image is delivered and the package declares no asset dependency.  Editorial formatting only, in addition: an AMS \texttt{amsart} preamble that restates the article's own declarations and macros used by the retained lines, namely the environments \texttt{lemma} on separate counters, together with the standard packages those lines need.  The retained environments print in this document as Lemma 1 in order of appearance, whereas the article prints the same items with the numbering of its own full document.  The complete native source of the article as distributed in the arXiv e-print for this exact version is bundled unchanged as \texttt{sources/201223/source0.tex}.
\begin{equation}
\ell_{\mathrm{Bayes}} \;\asymp\; \sum_{j>k} \operatorname{Var}(x_j).
\label{eq:bayes-tail}
\end{equation}

\begin{equation}
\ell_{\mathrm{Bayes},D}
\;\asymp\;
D^{-\frac{\alpha_1 + \alpha_2 + 1}{\alpha_1 + 1}},
\quad
\text{so that }
\alpha_D = \frac{\alpha_1 + \alpha_2 + 1}{\alpha_1 + 1}.
\label{eq:bayes-loss-over}
\end{equation}

\begin{lemma}[Learnable Coordinates in the Sparse Model when $\alpha_1>0$]
\label{lem:learnable-coordinates}
Under the sparse activation distribution
$p_j = \mathbb{P}(x_j \neq 0) = j^{-\alpha_1 - 1}$, the expected number of coordinates that are active in at least one of the $D$ training samples is asymptotic to the scale
\begin{equation}
K(D)
=
\Gamma\!\left(1 - \frac{1}{\alpha_1 + 1}\right)
D^{\frac{1}{\alpha_1 + 1}}, \footnote{When $\alpha_1>0$, $K(D)\asymp D^{1/(\alpha_1+1)}=o(D)$, \textit{i.e.} only a sublinear number of coordinates are ever observed. Interpolation is still possible since the data lie in a $K(D)$-dimensional subspace, which a linear predictor in $\mathbb{R}^N$ can fit. Unobserved coordinates receive no signal and are typically suppressed by mild regularization.
}
\label{eq:K-of-D}
\end{equation}
where $\Gamma(\cdot)$ is the Gamma function. Moreover, $K(D)$ \JSCheck{sets the scale of the} number of coefficients of $\mathbf{w}$ that can be estimated from data: a coordinate that is never activated carries no information about its associated weight, and hence cannot contribute to the learned predictor.
\end{lemma}

\begin{proof}[\JSCheck{Informal discrete argument via $K(D)$}]
By Lemma~\ref{lem:learnable-coordinates}, when $\alpha_1>0$, at most $K(D)$ coordinates of the
target vector $\mathbf{w}$ can be reliably estimated from $D$ samples. All
coordinates with indices $j > K(D)$ remain unobserved and therefore unmodeled.
By the unmodeled-variance principle (Eq.~\eqref{eq:bayes-tail}), the
Bayes-optimal loss is
\[
\ell_{\mathrm{Bayes},D}
\;\asymp\;
\sum_{j > K(D)} j^{-\alpha_1 - \alpha_2 - 2}.
\]
Approximating the sum by an integral,
\[
\begin{aligned}
\sum_{j > K(D)} j^{-\alpha_1 - \alpha_2 - 2}
&\;\asymp\;
\int_{K(D)}^\infty j^{-\alpha_1 - \alpha_2 - 2}\, dj
\\
&=
\frac{1}{\alpha_1 + \alpha_2 + 1}\,
K(D)^{-(\alpha_1+\alpha_2+1)}.
\end{aligned}
\]
Substituting the scaling $K(D) \asymp D^{1/(\alpha_1+1)}$ from
Eq.~\eqref{eq:K-of-D} yields
\[
\ell_{\mathrm{Bayes},D}
\;\asymp\;
D^{-(\alpha_1+\alpha_2+1)/(\alpha_1+1)},
\]
which is Eq.~\eqref{eq:bayes-loss-over}.
\end{proof}

\end{document}
